% ============================================================================
\section{Problems}
\label{sec:problems}
% ============================================================================

These are the problems that close the notebook.\ifsolutions{} Numerical answers quoted
in the solutions have been checked against direct diagonalization or propagation.\fi

% ---------------------------------------------------------------------------
\begin{problem}{The excitonic dimer}
\label{prob:exciton}
Two chromophores sit close enough that their transition dipoles interact. In the
singly-excited basis $\{\ket{e_1 g_2},\ket{g_1 e_2}\}$ the Hamiltonian is
\[
\Hop = \begin{bmatrix}\varepsilon_1 & J \\ J & \varepsilon_2\end{bmatrix},
\]
with $J$ the excitonic coupling.
\begin{enumerate}[label=(\alph*),leftmargin=*,itemsep=1pt]
\item Identify $\varepsilon_0$, $d_z$ and $d_x$.
\item For identical chromophores, find the eigenstates and energies. The splitting
$2|J|$ is the \emph{Davydov splitting}.
\item The transition dipole to the symmetric state is
$\bm{\mu}_1+\bm{\mu}_2$ and to the antisymmetric state $\bm{\mu}_1-\bm{\mu}_2$. In the
point-dipole approximation, $J>0$ for a side-by-side (H-aggregate) geometry and $J<0$
for head-to-tail (J-aggregate). In each case, is the \emph{bright} state the upper or
lower one, and does the absorption red- or blue-shift relative to the monomer?
\item Adapt the code from Part 3 of the notebook to plot the two exciton energies and
their oscillator strengths against $\varepsilon_1-\varepsilon_2$ at fixed $J$.
\end{enumerate}
\end{problem}

\begin{solution}
\textbf{(a)} By inspection, $\varepsilon_0 = \half(\varepsilon_1+\varepsilon_2)$,
$d_z = \varepsilon_1-\varepsilon_2$, $d_x = 2J$, $d_y = 0$.

\textbf{(b)} Identical chromophores means $d_z=0$, hence $\theta=\pi/2$ and the
eigenstates are $\ket{\pm} = \tfrac{1}{\sqrt2}\!\left(\ket{e_1g_2}\pm\ket{g_1e_2}\right)$
with $E_\pm = \varepsilon \pm J$. Note the sign: the symmetric combination is the
\emph{upper} state when $J>0$ and the \emph{lower} state when $J<0$.

\textbf{(c)} The oscillator strength goes as $|\bm{\mu}_1\pm\bm{\mu}_2|^2$. For two
parallel dipoles, the symmetric combination carries all of it and the antisymmetric
combination is dark. The point-dipole coupling is
$J = r^{-3}\!\left[\bm{\mu}_1\!\cdot\!\bm{\mu}_2 - 3(\bm{\mu}_1\!\cdot\!\hat{r})(\bm{\mu}_2\!\cdot\!\hat{r})\right]$:
side by side ($\bm{\mu}\perp\hat{r}$) gives $J=+\mu^2/r^3>0$, head to tail
($\bm{\mu}\parallel\hat{r}$) gives $J=-2\mu^2/r^3<0$.

\smallskip
\begin{center}\small
\begin{tabular}{@{}llll@{}}
\toprule
Geometry & Sign of $J$ & Bright (symmetric) state & Shift \\
\midrule
H-aggregate (side by side) & $J>0$ & upper, at $\varepsilon+J$ & \textbf{blue} \\
J-aggregate (head to tail) & $J<0$ & lower, at $\varepsilon+J$ & \textbf{red} \\
\bottomrule
\end{tabular}
\end{center}
\smallskip

\noindent
This is why J-aggregates show the narrow, red-shifted, intense band that makes them
useful as photographic sensitizers and light-harvesting mimics, while H-aggregates are
blue-shifted and comparatively dark.

\textbf{(d)} With $d_z\neq0$ the two chromophores are inequivalent, the states are no
longer exactly symmetric/antisymmetric, and the dark state steals intensity: the
oscillator strengths become $|\cos(\theta/2)\bm{\mu}_1 \pm \sin(\theta/2)\bm{\mu}_2|^2$
and neither state is completely dark. Energetic disorder therefore \emph{brightens}
H-aggregates --- a real and much-studied effect in molecular crystals.
\end{solution}

% ---------------------------------------------------------------------------
\begin{problem}{H\"uckel butadiene is two dimers}
\label{prob:butadiene}
The $\pi$ system of \emph{trans}-butadiene is a $4\times4$ H\"uckel problem
($H_{ii}=\alpha$, $H_{i,i+1}=\beta$), but it has a mirror symmetry exchanging
$1\leftrightarrow4$ and $2\leftrightarrow3$.
\begin{enumerate}[label=(\alph*),leftmargin=*,itemsep=1pt]
\item Form $s_{1,2}$ and $a_{1,2}$, the symmetric and antisymmetric combinations
$\tfrac{1}{\sqrt2}(p_1\pm p_4)$ and $\tfrac{1}{\sqrt2}(p_2\pm p_3)$, and show that
$\Hop$ block-diagonalizes into two $2\times2$ blocks.
\item Identify $d_z$ and $d_x$ for each block and use \cref{eq:eigenvalues} to obtain all
four $\pi$ energies. Check against the textbook result
$E = \alpha\pm1.618\beta,\ \alpha\pm0.618\beta$.
\item Verify by diagonalizing the full $4\times4$ matrix numerically.
\end{enumerate}
\end{problem}

\begin{solution}
\textbf{(a)} Sites 1 and 4 are not bonded, so $H_{14}=0$. Then
\[
\braket{s_1|\Hop|s_1} = \tfrac12(H_{11}+2H_{14}+H_{44}) = \alpha,
\qquad
\braket{s_2|\Hop|s_2} = \tfrac12(H_{22}+2H_{23}+H_{33}) = \alpha+\beta,
\]
\[
\braket{s_1|\Hop|s_2} = \tfrac12(H_{12}+H_{13}+H_{42}+H_{43}) = \beta,
\]
and every cross term $\braket{s_i|\Hop|a_j}$ vanishes because $\Hop$ commutes with the
mirror operation while $s$ and $a$ belong to different eigenvalues of it. Repeating for
the antisymmetric pair (the signs of $H_{23}$ and $H_{14}$ flip) gives
\[
\Hop_S = \begin{bmatrix}\alpha & \beta\\ \beta & \alpha+\beta\end{bmatrix},
\qquad
\Hop_A = \begin{bmatrix}\alpha & \beta\\ \beta & \alpha-\beta\end{bmatrix}.
\]

\textbf{(b)} Reading off the Pauli components:
\[
\begin{aligned}
\text{block } S:&\quad \varepsilon_0 = \alpha+\tfrac{\beta}{2}, &\quad d_z &= -\beta, &\quad d_x &= 2\beta;\\
\text{block } A:&\quad \varepsilon_0 = \alpha-\tfrac{\beta}{2}, &\quad d_z &= +\beta, &\quad d_x &= 2\beta.
\end{aligned}
\]
Both have $|\dvec| = \sqrt{\beta^2+4\beta^2} = \sqrt5\,|\beta|$, so
\[
E = \alpha \pm \frac{\beta}{2} \pm \frac{\sqrt5}{2}|\beta|
  \;=\; \alpha + 1.618\beta,\;\; \alpha+0.618\beta,\;\; \alpha-0.618\beta,\;\; \alpha-1.618\beta,
\]
since $\tfrac12(\sqrt5\pm1) = 1.618,\,0.618$ --- the golden ratio and its reciprocal,
which is where those famous H\"uckel numbers come from. With $\beta<0$ the two lowest
(occupied) levels are $\alpha+1.618\beta$ (from block $S$) and $\alpha+0.618\beta$
(from block $A$).

\textbf{(c)} Numerically, with $\alpha=0$ and $\beta=-1$: the full matrix gives
$\{-1.61803,\,-0.61803,\,+0.61803,\,+1.61803\}$, and the two blocks give
$\{-1.61803,\,+0.61803\}$ and $\{-0.61803,\,+1.61803\}$ respectively --- the same set.

\smallskip\noindent
\emph{Why this matters.} Reducing a large Hamiltonian to independent small blocks using
symmetry is one of the most useful manoeuvres in quantum chemistry, and here it turns a
$4\times4$ problem into two problems you can solve with \cref{eq:eigenvalues} in your
head. The same trick underlies symmetry-adapted linear combinations throughout group
theory.
\end{solution}

% ---------------------------------------------------------------------------
\begin{problem}{St\"uckelberg oscillations}
\label{prob:stuck}
Modify \texttt{landau\_zener\_numeric} so that instead of one linear sweep,
$d_z(t) = A\sin(\omega_{\rm sw}t)$ oscillates back and forth through the crossing.
\begin{enumerate}[label=(\alph*),leftmargin=*,itemsep=1pt]
\item Plot the population in $\ket{\alpha}$ against time over a few periods.
\item Explain why it does not simply relax toward $1/2$.
\item Find an amplitude and frequency for which the transfer after \emph{two} passages
is smaller than after one. What is interfering with what?
\end{enumerate}
\end{problem}

\begin{solution}
\textbf{(a),(b)} Each passage through the crossing acts as a beamsplitter, putting the
system into a superposition of the two adiabatic surfaces with amplitudes fixed by
\cref{eq:lz}. Between passages the two components accumulate \emph{different} dynamical
phases, $\Delta\phi = \hbar^{-1}\!\oint\!\left(E_+-E_-\right)\dd t$, because the surfaces
are separated. At the next passage the two amplitudes recombine and interfere. The
dynamics is therefore a repeated interferometer, not a repeated random draw --- so the
population does not diffuse to $1/2$, it oscillates. These are St\"uckelberg
oscillations, and the full sequence is a Landau--Zener--St\"uckelberg interferometer.

\textbf{(c)} Take $V_{DA}=0.35$, $A=6$ and scan $\omega_{\rm sw}$ with $\hbar=1$. The
ratio (transfer after two passages)/(transfer after one) is strongly non-monotonic ---
it runs from about $10$ down to about $1$ and back as $\omega_{\rm sw}$ varies over a
range of order unity. Two representative points:

\smallskip
\begin{center}\small
\begin{tabular}{@{}lccc@{}}
\toprule
$\omega_{\rm sw}$ & transfer after 1 passage & after 2 passages & ratio \\
\midrule
$0.450$ & $0.041$ & $0.430$ & $10.4$ \quad(constructive) \\
$0.700$ & $0.0244$ & $0.0241$ & $0.99$ \quad(destructive) \\
\bottomrule
\end{tabular}
\end{center}
\smallskip

\noindent
At $\omega_{\rm sw}\approx0.700$ the second passage almost exactly undoes the first: the
amplitude transferred on the way out and the amplitude transferred on the way back
arrive $\pi$ out of phase and cancel. Nothing has been ``untransferred'' in a classical
sense --- the two \emph{amplitudes} interfere destructively. Changing
$\omega_{\rm sw}$ by a few percent changes $\Delta\phi$ by order $\pi$ and flips the
result, which is why the ratio swings so violently.
\end{solution}

% ---------------------------------------------------------------------------
\begin{problem}{A genuine $\syo$: Peierls phases}
\label{prob:peierls}
A magnetic field attaches a phase to each hopping integral, $\beta\to\beta e^{\ii\phi_B}$.
\begin{enumerate}[label=(\alph*),leftmargin=*,itemsep=1pt]
\item Write $\Hop$ for a two-site dimer with this hopping and identify $d_x$ and $d_y$.
\item Show that the eigenvalues are independent of $\phi_B$, consistent with
\cref{sec:sigmay}.
\item Now take a triangle of three sites, each bond carrying phase $\phi_B$. Show that
the eigenvalues \emph{do} depend on the total flux $\Phi = 3\phi_B$ and that the
dependence cannot be rephased away. Diagonalize numerically and plot the three levels
against flux.
\end{enumerate}
\end{problem}

\begin{solution}
\textbf{(a)} With $H_{AB} = \beta e^{\ii\phi_B}$ and $\beta$ real,
$d_x = 2\beta\cos\phi_B$ and $d_y = -2\beta\sin\phi_B$, so
$\dperp = 2|\beta|$ regardless of $\phi_B$: the phase only rotates $\dvec$ within the
equatorial plane.

\textbf{(b)} By \cref{eq:eigenvalues} the eigenvalues depend on the transverse components
only through $\dperp$, so $E_\pm = \varepsilon_0 \pm |\beta|$ for every $\phi_B$.
Numerically, $\phi_B = 0,\,0.4,\,1.1,\,2.0$ all give $\{-|\beta|,+|\beta|\}$ exactly.
Equivalently: one rephasing, $\ket{B}\to e^{\ii\phi_B}\ket{B}$, removes the phase
completely.

\textbf{(c)} For the triangle with $t = \beta e^{\ii\phi_B}$ on each bond taken in a
consistent circulation,
\[
\Hop = \begin{bmatrix} 0 & t & t^{*}\\ t^{*} & 0 & t\\ t & t^{*} & 0\end{bmatrix},
\qquad
E_n = 2\beta\cos\!\left(\frac{2\pi n + \Phi}{3}\right),\quad n=0,1,2,
\qquad \Phi = 3\phi_B .
\]
This is verified numerically: at $\Phi=0$ the spectrum is
$\{-2|\beta|, +|\beta|, +|\beta|\}$ with the familiar doubly degenerate level of a
H\"uckel triangle; at $\Phi = \pi/2$ it is $\{-1.732,0,+1.732\}|\beta|$; at $\Phi = \pi$
it is $\{-|\beta|,-|\beta|,+2|\beta|\}$. The degeneracy is lifted and the whole spectrum
shifts with flux.

The counting argument is the point. Three sites admit three phase choices, but the
overall phase is unobservable, leaving \emph{two} useful rephasings against
\emph{three} bond phases. The combination that survives is the sum around the loop,
$\Phi = \phi_{12}+\phi_{23}+\phi_{31}$, which is gauge invariant and physically is the
magnetic flux threading the ring. A dimer has two sites and one bond --- one usable
rephasing, one phase --- so nothing survives. That is the whole difference between (b)
and (c), and it is the Aharonov--Bohm effect in its smallest possible setting.
\end{solution}

% ---------------------------------------------------------------------------
\begin{problem}{Ramsey interferometry}
\label{prob:ramsey}
Using the machinery of \cref{sec:driven}, simulate the following on a system with
detuning $\Delta$: a $\pi/2$ pulse with phase $\varphi=0$; free evolution with the drive
off (so $\dvec = \hbar\Delta\hat{z}$) for a time $\tau$; a second $\pi/2$ pulse, also
with $\varphi=0$.
\begin{enumerate}[label=(\alph*),leftmargin=*,itemsep=1pt]
\item Plot the final $P_\beta$ against $\tau$. What is the oscillation frequency?
\item Repeat with the second pulse phase-shifted to $\varphi=\pi$. What happens, and why
does this demonstrate that the \emph{relative} phase of two pulses is physical even
though the phase of one is not?
\item This is the Ramsey separated-oscillatory-fields method, and it is how atomic clocks
work. Why does a longer $\tau$ give a more precise frequency measurement?
\end{enumerate}
\end{problem}

\begin{solution}
Work in the hard-pulse limit $\omega_1\gg\Delta$, so the pulses are effectively
instantaneous rotations and the detuning acts only during $\tau$.

\textbf{(a)} The first $x$-pulse takes $\svec$ from the north pole to $-\hat{y}$
(\cref{sec:driven}). Free evolution rotates it in the equatorial plane at rate $\Delta$,
so $\svec(\tau) = (\sin\Delta\tau,\,-\cos\Delta\tau,\,0)$. The second $x$-pulse rotates
about $\hat{x}$ by $\ang{90}$ in the same sense ($\hat{y}\to\hat{z}$), leaving
$\svec = (\sin\Delta\tau,\,0,\,-\cos\Delta\tau)$. Hence
\[
P_\beta(\tau) = \frac{1-s_z}{2} = \frac{1+\cos\Delta\tau}{2},
\]
full-contrast fringes at the detuning frequency $\Delta$. (Numerically, with
$\omega_1 = 400\Delta$, the simulation matches this to better than $0.003$.)

\textbf{(b)} With $\varphi=\pi$ the second rotation is about $-\hat{x}$, so
$\hat{y}\to-\hat{z}$ and the final $s_z$ changes sign:
$P_\beta = \tfrac12(1-\cos\Delta\tau)$. The fringes are exactly inverted. Neither pulse
phase is meaningful alone --- calling the first one ``$x$'' was a convention --- but the
\emph{difference} of the two phases determines whether you measure a fringe maximum or a
minimum at $\tau=0$. That is a measurable, convention-independent statement, and it is
the concrete content of \cref{sec:sigmay}.

\textbf{(c)} The fringe pattern $\cos\Delta\tau$ has period $2\pi/\tau$ in $\Delta$, so a
frequency error $\delta\Delta$ produces a phase error $\delta\Delta\cdot\tau$. Doubling
$\tau$ halves the detuning uncertainty for the same phase resolution: the precision
improves as $1/\tau$. In practice $\tau$ is limited by decoherence, which is why
buying a better clock means buying a longer coherence time.
\end{solution}

% ---------------------------------------------------------------------------
\begin{problem}{The ammonia dipole paradox}
\label{prob:nh3}
At zero field the ground state of NH$_3$ is $\tfrac{1}{\sqrt2}(\ket{L}+\ket{R})$, which
has \emph{zero} expectation value of the dipole operator --- yet NH$_3$ is tabulated with
$\mu = \SI{1.47}{\debye}$.
\begin{enumerate}[label=(\alph*),leftmargin=*,itemsep=1pt]
\item Compute $\Braket{\hat\mu}$ in each of the two zero-field eigenstates, taking
$\hat\mu = \mu\,\szo$ in the $\{\ket{L},\ket{R}\}$ basis.
\item A molecule prepared in $\ket{L}$ is not stationary. Using \cref{eq:bloch-eq}, find
$\Braket{\hat\mu}(t)$ and its oscillation frequency.
\item Reconcile (a) and (b) with the tabulated value. On what timescale must a
measurement act to see $\SI{1.47}{\debye}$ rather than zero?
\item Deuterating to ND$_3$ drops the inversion splitting to about
\SI{0.053}{\wavenumber}. What happens to the crossover field
$\mathcal{E}_{\rm cross}=\Delta_t/\mu$, and why is the splitting so sensitive to
isotopic substitution?
\end{enumerate}
\end{problem}

\begin{solution}
\textbf{(a)} With $\hat\mu = \mu\szo$ and eigenstates
$\ket{\pm} = \tfrac{1}{\sqrt2}(\ket{L}\pm\ket{R})$,
\[
\Braket{\pm|\hat\mu|\pm} = \frac{\mu}{2}\big(\braket{L|\szo|L} + \braket{R|\szo|R}\big)
 = \frac{\mu}{2}(1-1) = 0 .
\]
Both stationary states are exactly nonpolar. This is not an accident: they are parity
eigenstates, and $\hat\mu$ is parity-odd, so its diagonal elements must vanish. The
dipole operator is purely \emph{off}-diagonal in the energy eigenbasis --- which is
precisely why the $\SI{23.786}{\giga\hertz}$ transition between them is dipole allowed
and drives the maser.

\textbf{(b)} $\ket{L} = \tfrac{1}{\sqrt2}(\ket{+}+\ket{-})$ has Bloch vector
$\svec(0) = +\hat{z}$, while $\dvec = -2\Delta_t\hat{x}$. By \cref{eq:bloch-eq} the
vector precesses in the $zy$ plane at $|\dvec|/\hbar = 2\Delta_t/\hbar$, so
\[
\Braket{\hat\mu}(t) = \mu\, s_z(t) = \mu\cos\!\left(\frac{2\Delta_t}{\hbar}t\right),
\]
oscillating at exactly the inversion frequency, \SI{23.786}{\giga\hertz}
(period \SI{42}{\pico\second}).

\textbf{(c)} Both are right; they answer different questions. A measurement slow
compared with \SI{42}{\pico\second} averages \cref{eq:bloch-eq} over many inversions and
returns zero. A measurement fast compared with it --- or any interaction that
\emph{localizes} the molecule before it can invert --- sees the instantaneous
\SI{1.47}{\debye}. In practice a bulk dielectric measurement reports
$\mu=\SI{1.47}{\debye}$ because collisions in the gas, and the local fields of
neighbouring molecules in a liquid, both act far faster than \SI{42}{\pico\second} and
keep the molecule pinned on one side. \Cref{fig:nh3} is the quantitative version of that
statement: a field of order $\mathcal{E}_{\rm cross}$ localizes the umbrella, and once
localized the molecule has its full classical dipole. The ``paradox'' is really the
statement that the tabulated dipole belongs to a \emph{localized} molecule, not to an
energy eigenstate of an isolated one.

\textbf{(d)} $\mathcal{E}_{\rm cross}$ is proportional to $\Delta_t$, so it falls by the
same factor of roughly $15$: from \SI{16.07}{\kilo\volt\per\centi\meter} to about
\SI{1.1}{\kilo\volt\per\centi\meter}. ND$_3$ is far easier to localize with a field.

The sensitivity comes from the exponential dependence of a tunneling amplitude on the
mass. In the WKB approximation
$\Delta_t \sim \exp\!\left(-\hbar^{-1}\!\int\!\sqrt{2m[V(x)-E]}\,\dd x\right)$, so
$\Delta_t$ falls roughly exponentially in $\sqrt{m}$. Deuterium is twice the mass of
hydrogen, and although the relevant inversion coordinate involves the nitrogen moving
against three hydrogens rather than a single H atom, the effective mass still increases
enough to suppress the amplitude by more than an order of magnitude. This exponential
mass sensitivity is the general signature of tunneling, and it is why kinetic isotope
effects far larger than the classical $\sqrt{2}$ are taken as evidence that a reaction
proceeds by tunneling.
\end{solution}
